% ============================================================================
%  Subdocumento 9. Contenido. Se usa igual en el documento único y en el
%  PDF propio del subdocumento. Los formularios que le asigna el T-21 se
%  agregan solos al final (datos en formularios/ de esta carpeta).
% ============================================================================

\begin{resumenApertura}
\marcador{Qué resuelve este subdocumento, en cuatro o cinco líneas}
\begin{recibe}
  \item \marcador{Qué recibe el mandante}
\end{recibe}
\end{resumenApertura}

\section{Marco de aseguramiento de calidad}\label{sec:9-marco-de-aseguramiento-de-calida}

\marcador{ISO/IEC 25010 y modelos de madurez}

\section{Métricas y umbrales bloqueantes}\label{sec:9-metricas-y-umbrales-bloqueantes}

\marcador{Calidad del código, cobertura, complejidad y acoplamiento}

\section{Puertas de calidad y revisiones}\label{sec:9-puertas-de-calidad-y-revisiones}

\marcador{Revisión por pares, análisis estático y dinámico}

\section{Estrategia de pruebas}\label{sec:9-estrategia-de-pruebas}

\marcador{ISO/IEC/IEEE 29119. El Formulario T-13 indica dónde está cada exigencia}

\section{Verificación, validación y trazabilidad}\label{sec:9-verificacion-validacion-y-trazab}

\marcador{De requerimiento a despliegue}
